\section{Introduction}
\label{sec:introduction}

The prevailing paradigm of large language model (LLM) capability improvement has
shifted from scaling training compute alone toward \emph{test-time compute scaling}:
investing additional inference-time resources to improve the quality of outputs for
individual queries \citep{snell2024scaling,brown2024large,openai2024o1}.
A central component of this paradigm is the \emph{verifier}---a model that scores
candidate partial or complete solutions---which guides a search procedure such as
best-of-$N$ sampling \citep{cobbe2021training}, beam search, or Monte Carlo Tree
Search (MCTS) \citep{silver2016mastering,guo2025deepseek}.

\paragraph{The efficiency problem.}
A key theoretical result of \citet{setlur2024rewarding} and the related analysis in
\citet{snell2024scaling} is that, among all test-time compute strategies, verifier-guided
search is asymptotically optimal: for a fixed compute budget, no method based on
unguided sampling can match the performance of a sufficiently accurate verifier.
Despite this optimality result, applying the verifier \emph{uniformly}---at every
node of the search tree or after every generated token---is computationally
prohibitive. Modern PRMs such as those of \citet{lightman2023lets} and
\citet{wang2024math} themselves carry significant inference cost; calling them
at every search step multiplies total compute by a large constant factor.

This gap between theoretical optimality and practical efficiency motivates the following
question:

\begin{quote}
\emph{Can we allocate verifier compute dynamically across search steps so that
total verification cost is sub-linear in the search budget, while suffering only
a bounded regret relative to exhaustive verification?}
\end{quote}

\paragraph{Our contributions.}
We answer this question affirmatively. Specifically:

\begin{enumerate}
  \item \textbf{Problem formalization.} We cast verifier allocation as an online
    resource allocation problem and define a precise regret criterion against an
    oracle that observes the optimal allocation in hindsight
    (Section~\ref{sec:formulation}).

  \item \textbf{Algorithm.} We introduce \textbf{DVA-MCTS} (Dynamic Verifier
    Allocation MCTS), which integrates a budget-sensitive uncertainty threshold
    into standard MCTS selection, determining at each node whether the verifier
    should be called (Section~\ref{sec:algorithm}).

  \item \textbf{Regret bound.} Under a Lipschitz score continuity assumption on
    the PRM (Assumption~\ref{ass:lipschitz}), we prove
    \[
      R(T) \;\leq\; C\sqrt{T \log K},
    \]
    where $T$ is the total number of search steps, $K$ is the branching factor,
    and $C$ is a constant depending on the Lipschitz parameter and verifier noise
    (Theorem~\ref{thm:main}). DVA-MCTS simultaneously invokes the verifier at most
    $N_{\mathcal{T}} = K(K^D-1)/(K-1)$ times total, so the call fraction $\mathcal{B}/T \to 0$
    as $T \to \infty$ with $K, D$ fixed. The threshold $\tau_t$ additionally
    provides a $T_{\mathrm{free}} = \lfloor(\gamma/(LD))^2\rfloor$-step call-free
    initial phase whose length grows with signal-to-noise ratio $\sigma_{\max}/L$
    (Theorem~\ref{thm:calls}).

  \item \textbf{Lower bound.} We prove that any online verifier allocation
    algorithm must incur $\Omega(\sqrt{T})$ regret on some problem instance,
    establishing that DVA-MCTS is rate-optimal (Theorem~\ref{thm:lower}).

  \item \textbf{Empirical validation.} Experiments with 50 matched-pair runs confirm
    the $O(\sqrt{T})$ regret rate (empirical slope $0.635$, $R^2 = 0.961$).
    A new exploitation-regime experiment ($B \in \{400, \ldots, 16{,}384\}$)
    directly validates Theorem~\ref{thm:calls}: DVA call counts grow by $2.7\times$
    while the budget grows $40\times$, achieving $\mathbf{93.7\%}$ call savings at
    $B{=}16{,}384$ with $100\%$ accuracy matching Uniform Verification.
    All pairwise comparisons in the exploration regime are statistically
    non-significant (Wilcoxon $p > 0.10$), confirming the theoretical prediction
    that large savings require $T \gtrsim N_{\mathcal{T}}$. An adaptive online estimator for $L$
    requires no advance knowledge of the Lipschitz constant and achieves the
    best empirical regret across all threshold configurations.  On $100$ real
    GSM8K problems (Math-Shepherd PRM), DVA-MCTS achieves $8.8\%$ fewer calls
    \emph{and} $3.6$ pp higher accuracy than exhaustive verification.
    With running-average continuous scores ($L_{\text{eff}}{=}0.155$), DVA-MCTS
    achieves $\mathbf{49.9\%}$ call savings, confirming that $L_{\text{eff}}$ is
    the primary driver of efficiency (Section~\ref{sec:experiments}).
\end{enumerate}

\paragraph{Significance.}
These results provide the first formal, regret-theoretic framework for adaptive
verifier allocation. Efficiency unfolds across three budget regimes:

\begin{itemize}
  \item \emph{Exploration} ($B \ll N_{\mathcal{T}}{=}8{,}190$ for $K{=}2$, $D{=}12$):
    DVA's threshold rarely fires on first-visit nodes; savings are modest (${<}5\%$).
    All algorithms are statistically indistinguishable ($p{>}0.10$, Wilcoxon).

  \item \emph{Transition} ($B \lesssim N_{\mathcal{T}}$):
    As the tree fills, the Lipschitz proxy substitutes for increasingly many verifier
    calls: $23\%$ savings at $B{=}800$, $51\%$ at $B{=}1{,}600$,
    $73\%$ at $B{=}3{,}200$ (requiring only $862$ calls vs.\ $3{,}200$ for Uniform).

  \item \emph{Exploitation} ($B \gg N_{\mathcal{T}}$):
    DVA call counts plateau (single-verification ceiling, Theorem~\ref{thm:calls}(a)):
    $\mathbf{93.7\%}$ savings at $B{=}16{,}384$ with $100\%$ accuracy,
    identical to Uniform Verification.
    A small-tree experiment ($K{=}2$, $D{=}4$, $N_{\mathcal{T}}{=}30$) makes this
    regime directly checkable at budget $B{=}400$: $94.6\%$ savings.
\end{itemize}

Crucially, \emph{no advance knowledge of $L$ is required}: the adaptive
$\hat{L}$ estimator tracks the running maximum of observed score-change ratios,
converges to $\hat{L}{=}0.431$ (a conservative $23\%$ over-estimate), and achieves
the lowest empirical regret ($23.88$) and $100\%$ accuracy.
The effective Lipschitz constant $L_{\text{eff}}$ is the primary efficiency driver.
On Math-Shepherd binary PRMs ($L_{\text{eff}}{=}0.5$), DVA achieves $8.8\%$ savings
\emph{and} $3.6$ pp accuracy gain. Converting to running-average continuous scores
($L_{\text{eff}}{=}0.155$) recovers $\mathbf{49.9\%}$ savings---confirming the
prediction that neural regression PRMs with $L_{\text{eff}} \lesssim 0.1$ will
exceed $73\%$ savings on production systems.

\paragraph{Organization.}
Section~\ref{sec:related} reviews related work. Section~\ref{sec:formulation}
states the formal problem. Section~\ref{sec:algorithm} presents DVA-MCTS.
Section~\ref{sec:theory} develops the theoretical analysis.
Section~\ref{sec:experiments} reports experiments. Section~\ref{sec:conclusion}
concludes.
